\documentclass{beamer}
\usepackage[T1]{fontenc}

% ---------------------------------------------------------------------------
% Paleta fiel ao material original desta aula (extraída por amostragem de
% pixels do PDF fonte): azul petróleo no cabeçalho/título, azul vivo nos
% blocos, vermelho nos marcadores e amarelo/laranja/verde nos arrays.
% ---------------------------------------------------------------------------
\definecolor{hheader}{RGB}{0,102,153}
\definecolor{hblock}{RGB}{1,146,225}
\definecolor{hblockbg}{RGB}{191,228,248}
\definecolor{haccent}{RGB}{215,15,3}
\definecolor{hband}{RGB}{225,225,225}
\definecolor{hyellow}{RGB}{255,204,0}
\definecolor{horange}{RGB}{255,102,0}
\definecolor{hgreen}{RGB}{159,212,0}
\definecolor{hcode}{RGB}{60,60,60}

\setbeamercolor{title in head/foot}{bg=hheader,fg=white}
\setbeamercolor{section in head/foot}{bg=hband,fg=haccent}
\setbeamercolor{subsection in head/foot}{bg=hband,fg=haccent}
\setbeamercolor{frametitle}{bg=hheader,fg=white}
\setbeamerfont{frametitle}{series=\bfseries}
\setbeamercolor{block title}{bg=hblock,fg=white}
\setbeamercolor{block body}{bg=hblockbg,fg=black}
\setbeamercolor{itemize item}{fg=haccent}
\setbeamercolor{itemize subitem}{fg=haccent}
\setbeamercolor{itemize subsubitem}{fg=haccent}
\setbeamercolor{section in toc}{fg=haccent}
\setbeamercolor{subsection in toc}{fg=black!75}
\setbeamertemplate{navigation symbols}{}
\setbeamertemplate{itemize items}[triangle]
\setbeamertemplate{caption}[numbered]

% Cabeçalho: mesmo template do outer theme "tree" (usado pelo Montpellier),
% com a contagem de página adicionada na barra superior.
\makeatletter
\setbeamertemplate{headline}{%
    \begin{beamercolorbox}[wd=\paperwidth,colsep=1.5pt]{upper separation line head}
    \end{beamercolorbox}
    {\usebeamerfont{title in head/foot}%
    \begin{beamercolorbox}[wd=\paperwidth,ht=2.5ex,dp=1.125ex,leftskip=.3cm,rightskip=.3cm plus1fil]{title in head/foot}
      \insertshorttitle\hfill\insertframenumber{} / \inserttotalframenumber
    \end{beamercolorbox}}
    {\usebeamerfont{section in head/foot}%
    \begin{beamercolorbox}[wd=\paperwidth,ht=2.5ex,dp=1.125ex,leftskip=.3cm,rightskip=.3cm plus1fil]{section in head/foot}
        \setbox\beamer@tempbox=\hbox{\insertsectionhead}%
        \ifdim\wd\beamer@tempbox>1pt%
          \hskip2pt\raise1.9pt\hbox{\vrule width0.4pt height1.875ex\vrule width 5pt height0.4pt}\hskip1pt%
        \fi%
      \insertsectionhead
    \end{beamercolorbox}}
    {\usebeamerfont{subsection in head/foot}%
    \begin{beamercolorbox}[wd=\paperwidth,ht=2.5ex,dp=1.125ex,leftskip=.3cm,rightskip=.3cm plus1fil]{subsection in head/foot}
        \setbox\beamer@tempbox=\hbox{\insertsubsectionhead}%
        \ifdim\wd\beamer@tempbox>1pt%
          \hskip9.4pt\raise1.9pt\hbox{\vrule width0.4pt height1.875ex\vrule width 5pt height0.4pt}\hskip1pt%
        \fi%
      \insertsubsectionhead
    \end{beamercolorbox}}
    \begin{beamercolorbox}[wd=\paperwidth,colsep=1.5pt]{lower separation line head}
    \end{beamercolorbox}
}


% listings precisa estar carregado antes de \lstdefinestyle abaixo.
\usepackage{listings}

% Definição do estilo para Python
\lstdefinestyle{PythonStyle}{
    language=Python,
    basicstyle=\ttfamily,
    keywordstyle=\bfseries\color{blue},
    commentstyle=\itshape\color{green!60!black},
    stringstyle=\color{orange},
    showstringspaces=false,
    columns=flexible,
    frame=single,
    breaklines=true,
    literate={None}{{\textbf{None}}}4,
    morekeywords={if, class,for, self},
}

% Definição do estilo para C
\lstdefinestyle{CStyle}{
    language=C,
    basicstyle=\ttfamily,
    keywordstyle=\bfseries\color{blue},
    commentstyle=\itshape\color{green!60!black},
    stringstyle=\color{orange},
    showstringspaces=false,
    columns=flexible,
    frame=single,
    breaklines=true,
    literate={NULL}{{\textbf{NULL}}}4,
}


\setbeamerfont{headline}{size=\footnotesize}

\mode<presentation>{
\usetheme{Montpellier}
\setbeamercovered{transparent}
\usecolortheme{lsc}
}

\mode<handout>{
  % tema simples para ser impresso
  \usepackage[bar]{beamerthemetree}
   % Colocando um fundo cinza quando for gerar transparências para serem impressas
  % mais de uma transparência por página
  \beamertemplatesolidbackgroundcolor{black!5}
}


\usepackage{amsmath,amssymb}
\usepackage[brazil]{varioref}
\usepackage[english,brazil]{babel}
\usepackage[utf8]{inputenc}
\usepackage{fontawesome}

%\usepackage[latin1]{inputenc}
\usepackage{graphicx}
\usepackage{listings}
\usepackage{url}
\usepackage{colortbl}
\usepackage{caption}
\usepackage{tikz}
\usetikzlibrary{arrows.meta,positioning,shapes.geometric,calc,fit}
\usepackage{transparent}
\usepackage{multirow}
\usepackage{graphicx}
\usepackage{listings}
\usepackage{exercise,chngcntr}

\definecolor{maroon}{cmyk}{0, 0.87, 0.68, 0.32}
\definecolor{halfgray}{gray}{0.55}
\definecolor{ipython_frame}{RGB}{207, 207, 207}
\definecolor{ipython_bg}{RGB}{247, 247, 247}
\definecolor{ipython_red}{RGB}{186, 33, 33}
\definecolor{ipython_green}{RGB}{0, 128, 0}
\definecolor{ipython_cyan}{RGB}{64, 128, 128}
\definecolor{ipython_purple}{RGB}{170, 34, 255}

\beamertemplatetransparentcovereddynamic
\addtobeamertemplate{title page}{\centering{ \includegraphics[width=1.25cm]{figs/logo.png}}}{}
\title[Estrutura de Dados II]{Hash}
\author[Elton Sarmanho]{ 
Professor: Elton Sarmanho\inst{1} 
\newline
E-mail: \href{mailto:eltonss@ufpa.br}{eltonss@ufpa.br}\newline
\href{https://github.com/eltonsarmanho}{\faGithub} \href{https://www.linkedin.com/in/elton-sarmanho-836553185/}{\faLinkedinSquare}
}   
  \institute[UFPA]{
   \inst{1}%
     Faculdade de Sistemas de Informação - UFPA/CUTINS
     }     
 
% Se comentar a linha abaixo, irá aparecer a data quando foi compilada a apresentação  
\date{\today}
\pgfdeclareimage[height= 1.25 cm]{inf}{figs/logo.png}
% pode-se colocar o LOGO assim
\logo{{\transparent{0.35}\pgfuseimage{inf}}}

% \AtBeginSection[]{
%   \begin{frame}<beamer>
%     \frametitle{Roteiro}
%     \tableofcontents[currentsection,currentsubsection]
%   \end{frame}  
% }

\setbeamertemplate{navigation symbols}{\insertframenavigationsymbol \insertsectionnavigationsymbol \insertbackfindforwardnavigationsymbol} 


\expandafter\def\expandafter\insertshorttitle\expandafter{%
  \insertshorttitle\hfill%
  \insertframenumber\,/\,\inserttotalframenumber}
  \renewcounter{Exercise}[section]
  \newcounter{Example}

% ---------------------------------------------------------------------------
% Definições usadas pelos diagramas nativos em TikZ da seção "Hashing na
% Prática" (cores e estilos de nó/seta que faltavam neste arquivo e faziam
% os diagramas compilar sem cor/borda/seta).
% ---------------------------------------------------------------------------
\definecolor{hyellow}{RGB}{255,204,0}
\definecolor{horange}{RGB}{255,102,0}
\definecolor{hgreen}{RGB}{159,212,0}
\definecolor{hcode}{RGB}{60,60,60}
\tikzset{
  cell/.style={draw=black!70,thick,minimum width=6.5mm,minimum height=6.5mm,inner sep=0pt,font=\scriptsize,anchor=center},
  cellY/.style={cell,fill=hyellow},
  cellO/.style={cell,fill=horange},
  cellG/.style={cell,fill=hgreen},
  cellW/.style={cell,fill=white},
  idx/.style={font=\tiny,text=black!55},
  slot/.style={draw=black!70,thick,minimum width=9mm,minimum height=6mm,inner sep=1pt,font=\scriptsize},
  keytag/.style={draw=hheader,thick,rounded corners=1pt,fill=hblockbg,font=\scriptsize,inner sep=1.5pt,minimum width=7mm},
  uset/.style={ellipse,draw=black!70,thick,minimum width=26mm,minimum height=20mm},
  harrow/.style={-{Stealth[length=1.6mm]},thick,draw=black!70},
  redarrow/.style={-{Stealth[length=1.6mm]},thick,draw=haccent},
  chainnode/.style={draw=hheader,thick,fill=hblockbg,minimum width=5.5mm,minimum height=5.5mm,inner sep=0pt,font=\tiny},
}

\begin{document}

\setbeamertemplate{logo}{}
\begin{frame}
\titlepage
\end{frame}
%Começar mostrar Logo apos slide de titulo
\logo{{\transparent{0.35}\pgfuseimage{inf}}}

\begin{frame}
\frametitle{Roteiro}
\tableofcontents[sections={1-2}]
\end{frame}
 \begin{frame}
 \frametitle{Roteiro}
 \tableofcontents[sections={3-4}]
 \end{frame}


%%%%%%%%%%%%%%%%%%%%%%%%%%%%%%%%%%%%%%%%
\section{Hash}

\subsection{Introdução}
\begin{frame}
 \begin{itemize}
  \item Os métodos de busca linear ou binária baseiam-se na comparação da chave de busca com as chaves dos registros armazenados na tabela.
 \end{itemize}
\end{frame}

\begin{frame}
  \begin{columns}
 		  \column{0.39\linewidth}
 		  \small
          \begin{block}{Busca Sequencial}
 		  \begin{description}
              \item[Melhor Caso] $O(1)$ 
              \item[Pior Caso] $O(n)$ 
             \end{description} 
 		  \end{block}           
          \column{0.7\linewidth}
        \begin{figure}[h]
        \centering 
        \includegraphics[bb=0 0 141 74,scale=1.3]{figs/buscasequencial.png}
        %\caption*{Nomenclatura dos Componentes}
        \end{figure}          
  \end{columns}
\end{frame}

\begin{frame}
  \begin{columns}
 		  \column{0.39\linewidth}
 		  \small
          \begin{block}{Busca binária}
          Chave é comparada com registro que se encontra no meio do array ordenado
 		  \begin{description}
               
              \item[Melhor Caso] $O(1)$ 
              \item[Pior Caso] $O(\log n)$ 
             \end{description} 
 		  \end{block}           
          \column{0.7\linewidth}
        \begin{figure}[h]
        \centering 
        \includegraphics[bb=0 0 75 62,scale=1.8]{figs/buscaBinaria.png}

        %\caption*{Nomenclatura dos Componentes}
        \end{figure}          
  \end{columns}
\end{frame}

\begin{frame}
 \begin{itemize}
  \item A busca pode ser mais rápida se existir uma relação direta entre a posição em que o registro se encontra na tabela e a sua chave a ser localizada.
  \item Organizamos os dados em uma nova estrutura de dados chamada \textbf{Tabela Hash}
  \begin{itemize}
  \item Média por operação de $\theta(1)$, sendo o pior caso, entretanto, $O(n)$
  \end{itemize}

  
 \end{itemize}
\end{frame}

\subsection{Características}
\begin{frame}
 \begin{itemize}
  \item Os registros são armazenados em uma tabela $T$, sequencial e de dimensão $m$.
  \item As posições da tabela que se situam no intervalo $[0, m – 1]$ são calculadas através de sua chave por meio de uma função de dispersão (\textit{Hash Function})
  \begin{itemize}
   \item $h: U \rightarrow \{0,1,...,m-1\} $
   \item Universo das chaves $U$
   \item Conjunto dos endereços $\{0,1,...,m-1\}$ 
  \end{itemize}
  \item $h$ mapeia o universo $U$ de chaves nos endereços de uma tabela $T$ chamada de \textbf{tabela hash}(\textit{hash table})
 \end{itemize}
\end{frame}

\begin{frame}
 \begin{columns}
 		  \column{0.39\linewidth}
 		  \small

          \begin{block}{Técnica de Acesso Direto} 		   
            \begin{itemize}
            \item número de chaves $n$
            \item número de endereços da tabela $m$
            \item $n\leq m$
            \item Todas operações $O(1)$
            \end{itemize} 
 		  \end{block}           
          \column{0.7\linewidth}
        \begin{figure}[h]
        \centering
       \includegraphics[bb=0 0 231 113,scale=0.8]{figs/acesso_direto.png}

        \end{figure}      
           
  \end{columns}

 \end{frame}
 
\subsection{Problemas}
\begin{frame}
 \begin{itemize}
  \item Desperdício de espaço.
  \begin{itemize}
   \item Exemplo: Se existem duas chaves com índices $0$ e $999.999$
   \item A aplicação da técnica de acesso direto conduziria a uma tabela com $1.000.000$ compartimentos, dos quais apenas dois ocupados
  \end{itemize}
  \item Quando $n \geq m$
  \item \textbf{Solução}: aplicar $h(k)$ $\forall$ $k \in U$
  \begin{itemize}
   \item A ideia é transformar cada chave $k$ em um valor no intervalo $[0, m – 1]$ usando $h(k)$
   \item Dizemos que $h(k)$ é \textit{hash value} da chave $k$ 
  \end{itemize}  
 \end{itemize}
\end{frame}



\begin{frame}

\begin{columns}
 		  \column{0.5\linewidth}
 		   \begin{figure}[h]
 \centering
 \includegraphics[bb=0 0 281 100,scale=1.]{figs/funcaoHash.png}
 \caption*{Função Hash}
\end{figure}          
          \column{0.6\linewidth}
       \begin{figure}[h]
 \centering
 \includegraphics[bb=0 0 170 110,scale=1]{figs/acesso_direto_2.png}
 \caption*{Acesso Direto}
\end{figure}
           
  \end{columns}
\end{frame}

\begin{frame}
 \begin{itemize}
  \item Função de dispersão pode não garantir \textbf{injetividade}
 \begin{itemize}
\item É possível a existência de outra chave $y \neq x$ 
\item Duas chaves podem ter mesmo \textbf{valor hash} para o mesmo endereço 
\begin{itemize}
 \item $h(x) = h(y)$
\end{itemize}

\item Fenômeno chamado de \textbf{Colisão}
\end{itemize}   
 \end{itemize}
\end{frame}

\begin{frame}
\begin{figure}[h]
 \centering
 \includegraphics[bb=0 0 104 79,scale=1.4]{figs/colisao.png}
 \caption*{Exemplo de Colisão usando $h(k) = k$ $mod$ $5$}
\end{figure}

\end{frame}


\section{Funções de Dispersão}
\subsection{Propriedades}
\begin{frame}
 \begin{itemize}
  \item Uma função de dispersão $h$ transforma uma chave $x$ em um endereço-base $h(x)$ em $T$
  \item Uma $h(x)$ deve satisfazer às seguintes condições:
  \begin{itemize}
   \item produzir um número baixo de colisões;
   \begin{itemize}
    \item Tratamento de colisões
   \end{itemize}
   \item ser facilmente computável;   
   \item ser uniforme.
   \begin{itemize}
    \item A probabilidade de que $h(x)$ seja igual ao endereço $k$ deve ser $\frac{1}{m}$ para todas as chaves $x$ e todos os endereços $k \in [0,m-1]$
   \end{itemize}
  \end{itemize}
 \end{itemize}
\end{frame}

\subsection{Tipos de Funções}

\begin{frame}{Método da Divisão}
 \begin{itemize}
  \item Este método é fácil e eficiente, sendo por isso muito empregado. 
 \item A chave $x$ é dividida pela dimensão $m$ da tabela $T$, e o resto da divisão é usado como endereço.
 $$h(x) = x  \ \mathbf{mod}\ m$$ 
 \item Resultando em endereços no intervalo $[0,m-1]$
 \end{itemize} 
\end{frame}

\begin{frame}{Método da Divisão}
 \begin{itemize}
  \item Ao usar o método da divisão, geralmente evitamos certos valores de $m$. Existem critérios para escolher um melhor $m$.
 \begin{itemize}
  \item Escolher $m$ de modo que seja um número primo não próximo a uma potência de 2.  
 \end{itemize}
 \end{itemize} 
\end{frame}

\begin{frame}{Método da Divisão}
 \begin{figure}[h]
 \centering
 \includegraphics[bb=0 0 99 155,scale=.9]{figs/metodoDivisao.png}
 % metodoDivisao.png: 414x644 px, 300dpi, 3.51x5.45 cm, bb=0 0 99 155
 \caption*{Método da Divisão.\\ $\mathbf{m=23}$\\ $44\ \mathbf{mod} \ 23 = 21$}
\end{figure}

\end{frame}




\begin{frame}{Método da Multiplicação}
\begin{columns}
 		  \column{0.4\linewidth}
 		  \footnotesize
          \begin{block}{} 		   
            \begin{itemize}
  \item Endereço é calculado em duas etapas.
  \item A chave $k$ é multiplicada por uma constante $A \ \in \ ]0,1[$. O resultado é armazenado em $s$. Extrai a parte fracionada de $kA$
  \item Multiplica parte fracionada por $m$ e computa \textit{floor} $\lfloor \rfloor$ .
  $$h(k) = \lfloor m \ (kA \ \mathbf{mod} \ 1)\rfloor$$  
 \end{itemize} 
 		  \end{block}           
          \column{0.7\linewidth}
        \begin{figure}[h]
        \centering       
 \includegraphics[bb=0 0 122 64,scale=1.3]{figs/metodoMulti.png}
 % metodoMulti.png: 508x267 px, 300dpi, 4.30x2.26 cm, bb=0 0 122 64
        \end{figure}      
           
  \end{columns}

\end{frame}

\begin{frame}{Método da Multiplicação}
\begin{itemize}
 \item Vantagem desse método é que valor de $m$ não é crítico
 \item Embora esse método funcione com qualquer valor da constante $A$, ele funciona melhor com alguns valores do que com outros. 
 \begin{itemize}
\item A escolha ideal depende das características dos dados que estão sendo armazenados 
\item Knuth sugere:
$$A \approx \ \dfrac{(\sqrt{5}-1)}{2} = 0.6180...$$
\end{itemize}

\end{itemize}


\end{frame}

\begin{frame}{Exercício A - Método de Divisão}
\begin{itemize}
 \item Considere uma Tabela hash de tamanho $m = 8$. Calcule as localizações para as quais são mapeadas as chaves $[16, 23, 41, 25]$.

\end{itemize}
\end{frame}


\begin{frame}{Exercício A.1 - Método de Divisão}
\begin{itemize}
 \item Desenvolva a tabela hash com 9 slots inserindo $[$5, 28, 19, 15, 20, 33, 12, 17, 10$]$.
\end{itemize}
\end{frame}

\begin{frame}[fragile]{Implementação em C - Método da Divisão}
\footnotesize
\begin{lstlisting}[style=CStyle]
#include <stdio.h>
#define TABLE_SIZE 8
#include <math.h>
void initialize_table(int table[], int size) {
    for (int i = 0; i < size; i++) {
        table[i] = -1;// Valor inicial para indicar posicao vazia
    }
}

void insere(int table[], int key, char funcao) {
    if(funcao == 'D'){//Metodo da Divisao
    int hash = key % TABLE_SIZE;
    table[hash] = key;
    }
}

\end{lstlisting}
\end{frame}


\begin{frame}[fragile]{Implementação em C - Método da Divisão}
\footnotesize
\begin{lstlisting}[style=CStyle]
//Continuacao
int main() {
    int hash_table[TABLE_SIZE];
    initialize_table(hash_table, TABLE_SIZE);
    int keys[] = {16,23,41,25};
    int num_keys = sizeof(keys) / sizeof(keys[0]);

    for (int i = 0; i < num_keys; i++) {
        int key = keys[i];
        insere(hash_table, key,'D');
    }
    // Imprimir tabela hash
    for (int i = 0; i < TABLE_SIZE; i++) {
        printf("Index: %d, Key: %d\n", i, hash_table[i]);
    }
    return 0;
}
\end{lstlisting}
\end{frame}


\begin{frame}[fragile]{Implementação em Python - Método da Divisão}
 
\centering\footnotesize
\begin{lstlisting}[style=PythonStyle]
import math
import numpy as np

class Hashing:
    def __init__(self,funcao):
        self.funcao = funcao
        pass;
    
\end{lstlisting}


\end{frame} 

\begin{frame}[fragile]{Implementação em Python - Método da Divisão}
\centering\footnotesize
\begin{lstlisting}[style=PythonStyle]
    def hashingDivisao(self, chave, hashTable):
        m = len(hashTable)
        return chave % m;
    # Funcao add valor na Tabela hash
    def insere(self,hashtable, valor_chave, valor):
        if(self.funcao == "Divisao"):
            endereco = self.hashingDivisao(valor_chave,hashtable)
        
        hashtable[endereco]= valor
    
\end{lstlisting}


\end{frame} 

\begin{frame}[fragile]{}
\footnotesize
\begin{lstlisting}[style=PythonStyle]
if __name__ == '__main__':
    HashTable = np.array([None]*8)
    hashing = Hashing(funcao="Divisao")
    hashing.insere(HashTable, 16, '16')
    hashing.insere(HashTable, 23, '23')
    hashing.insere(HashTable, 41, '41')
    hashing.insere(HashTable, 25, '25')
    hashing.display_hash(HashTable)    
\end{lstlisting}
\footnotesize
\begin{block}{Resultado}
$0 --> 16$ \\
$1 --> 25$ \\
$2 --> None$\\
$3 --> None$\\
$4 --> None$\\
$5 --> None$\\
$6 --> None$\\
$7 --> 23$\\
\end{block}

\end{frame} 

\begin{frame}{Exercício B - Método da Multiplicação}
\begin{itemize}
 \item Considere uma Tabela hash de tamanho $m = 1000$ e uma função hash correspondente $h(k)$ igual a $m \ (k A\ mod \ 1)$ para $A \approx \ \dfrac{(\sqrt{5}-1)}{2}$. Calcule as localizações para as quais são mapeadas as chaves $61, 62, 63, 64 \ e \ 65$.

\end{itemize}
\end{frame}

\begin{frame}{Exercício B - Método da Multiplicação}
\begin{itemize}
 \item m =1000
 \item k = [61]
 \item $A = \ \dfrac{(\sqrt{5}-1)}{2}$ 
 $$h(k) = \lfloor m \ (kA \ \mathbf{mod} \ 1)\rfloor$$  
\end{itemize}
\end{frame}

\begin{frame}{Exercício B - Método da Multiplicação}
\begin{itemize}
 \item m =1000
 \item k = [61]
 \item $A = \ \dfrac{(\sqrt{5}-1)}{2}$ 
 $$h(k) = \lfloor m \ (k\dfrac{(\sqrt{5}-1)}{2} \ \mathbf{mod} \ 1)\rfloor$$  
\end{itemize}
\end{frame}

\begin{frame}{Exercício B - Método da Multiplicação}
\begin{itemize}
 \item m =1000
 \item k = [61]
 \item $A = \ \dfrac{(\sqrt{5}-1)}{2}$ 
 $$h(61) = \lfloor m \ (61\dfrac{(\sqrt{5}-1)}{2} \ \mathbf{mod} \ 1)\rfloor$$  
\end{itemize}
\end{frame}

\begin{frame}{Exercício B - Método da Multiplicação}
\begin{itemize}
 \item m =1000
 \item k = [61]
 \item $A = \ \dfrac{(\sqrt{5}-1)}{2}$ 
 $$h(61) = \lfloor m \ (37,700073314 \ \mathbf{mod} \ 1)\rfloor$$  
\end{itemize}
\end{frame}

\begin{frame}{Exercício B - Método da Multiplicação}
\begin{itemize}
 \item m =1000
 \item k = [61]
 \item $A = \ \dfrac{(\sqrt{5}-1)}{2}$ 
 $$h(61) = \lfloor m \ (0,700073314)\rfloor$$  
\end{itemize}
\end{frame}

\begin{frame}{Exercício B - Método da Multiplicação}
\begin{itemize}
 \item m =1000
 \item k = [61]
 \item $A = \ \dfrac{(\sqrt{5}-1)}{2}$ 
 $$h(61) = \lfloor 1000 \ (0,700073314)\rfloor$$  
\end{itemize}
\end{frame}

\begin{frame}{Exercício B - Método da Multiplicação}
\begin{itemize}
 \item m =1000
 \item k = [61]
 \item $A = \ \dfrac{(\sqrt{5}-1)}{2}$ 
 $$h(61) = \lfloor 700,073314\rfloor$$  
\end{itemize}
\end{frame}

\begin{frame}{Exercício B - Método da Multiplicação}
\begin{itemize}
 \item m =1000
 \item k = [61]
 \item $A = \ \dfrac{(\sqrt{5}-1)}{2}$ 
 $$h(61) = \lfloor 700,073314\rfloor$$  
\end{itemize}
\end{frame}

\begin{frame}{Exercício B - Método da Multiplicação}

\begin{itemize}
 \item m =1000
 \item k = [61]
 \item $A = \ \dfrac{(\sqrt{5}-1)}{2}$ 
 $$h(61) = 700$$  
\end{itemize}
\end{frame}

\begin{frame}[fragile]{Implementação em C - Método da Multiplicação}
\footnotesize
\begin{lstlisting}[style=CStyle]
//Alteramos somente metodo insere
void insere(int table[], int key, char funcao) {
    if(funcao == 'D'){
        int hash = key % TABLE_SIZE;
        table[hash] = key;
    }
    else if(funcao == 'M')
    {
        float A = (sqrt(5)-1)/2;
        float m = TABLE_SIZE;

        int hash = floor(m*(fmod(key*A , 1)));
        table[hash] = key;
    }
}

\end{lstlisting}
\end{frame}

\begin{frame}[fragile]{Implementação em C - Método da Multiplicação}
\footnotesize
\begin{lstlisting}[style=CStyle]
//Caso esteja usando linux sera preciso implementar
//metodo floor e fmod (biblioteca math.h)
double floor(double x) {
    int int_part = (int)x;
    return (double)int_part;
}
double fmod(double x, double y) {
    if (y == 0.0) {
        // Tratar divisao por zero
        return 0.0;
    }

    double quotient = x / y;
    double whole_part = (double)((int)quotient);
    double remainder = x - (whole_part * y);
    return remainder;
}

\end{lstlisting}
\end{frame}


\begin{frame}[fragile]{Implementação em C - Método da Multiplicação}
\footnotesize
\begin{lstlisting}[style=CStyle]
#define TABLE_SIZE 1000
.
int main() {
    int hash_table[TABLE_SIZE];
    initialize_table(hash_table, TABLE_SIZE);
    int keys[] = {61,62,63,64,65};
    int num_keys = sizeof(keys) / sizeof(keys[0]);
    for (int i = 0; i < num_keys; i++) {
        int key = keys[i];
        insere(hash_table, key,'M');
    }
    for (int i = 0; i < TABLE_SIZE; i++) {    // Imprimir tabela
        printf("Index: %d, Key: %d\n", i, hash_table[i]);
    }
    printf("Index: %d, Key: %d\n", 700, hash_table[700]);
    return 0;
}

\end{lstlisting}
\end{frame}



\begin{frame}[fragile]{Implementação em Python - Método da Multiplicação}
\centering\footnotesize
\begin{lstlisting}[style=PythonStyle]
    def hashingMultiplicacao(self,chave,hashTable):
        A = (math.sqrt(5)-1)/2
        m = len(hashTable)
        h = math.floor(m*(chave*A % 1))
        return h;

    def insere(self,hashtable, valor_chave, valor):
        if(self.funcao == "Divisao"):
            endereco = self.hashingDivisao(valor_chave,hashtable)
        elif (self.funcao == "Multiplicacao"):
            endereco = self.hashingMultiplicacao(valor_chave,hashtable)
        hashtable[endereco]=valor
    
\end{lstlisting}


\end{frame} 

\begin{frame}[fragile]{Implementação em Python - Método da Multiplicação}
\centering\footnotesize
\begin{lstlisting}[style=PythonStyle]
if __name__ == '__main__':
    HashTable = np.array([None] * 1000)
    hashing = Hashing(funcao="Multiplicacao")
    # Add os elementos na Tabela
    hashing.insere(HashTable, 61, '61',)
    hashing.insere(HashTable, 62, '62',)
    hashing.insere(HashTable, 63, '63', )
    hashing.insere(HashTable, 64, '64')
    hashing.insere(HashTable, 65, '65')
    hashing.display_hash(HashTable)   
\end{lstlisting}
\begin{block}{Resultado}
$700 --> 61$ \\
\end{block}

\end{frame} 

%
% \subsection{Trabalho I - Função de dispersão}
% \begin{frame}
% \begin{itemize}
%  \item Implemente o \textbf{Método da dobra}:
%  \begin{itemize}
%  \item Codifique usando chaves numéricas
%  \item Codifique usando chaves de caracteres
%  \item Não é necessário implementar Tratamento de colisões
%  \end{itemize}
%
% \end{itemize}
% \end{frame}
% \subsection{Trabalho II - Função de dispersão}
% \begin{frame}
% \begin{itemize}
%  \item Faça uma lauda sobre a relação de \textit{hashing} e \textit{blockchain}
%  \begin{itemize}
%   \item N equipes serão sorteadas aleatoriamente para apresentar esse trabalho.
%   \item Faça slides
%  \end{itemize}
%
% \end{itemize}
% \end{frame}

\section{Tratamento de Colisões}
\subsection{Fator de Carga}
\begin{frame}
 \begin{itemize}
  \item Já foi observado que o mesmo $h(x) = h(y)$ endereço pode ser encontrado para chaves diferentes $x \neq y$
  \begin{itemize}
   \item Colisão
  \end{itemize}
  \item Dado uma tabela $\mathbf{T}$ com $\mathbf{m}$ endereços e que armazena $\mathbf{n}$ elementos definimos \textbf{fator de carga} (\textit{load factor}):
  \begin{itemize}
   \item \emph{número médio de elementos armazenados em uma entrada}
   \item $\alpha = \dfrac{n}{m}$, onde  $\alpha < 1   $, $ \alpha = 1$ e $ \alpha > 1$
   \item $\uparrow \alpha \sim \ \uparrow $ Colisões 
  \end{itemize}
  \item \textbf{Encadeamento}
 \end{itemize}

\end{frame}

\subsection{Encadeamento}
\begin{frame}

\begin{columns}
        \column{0.3\linewidth}
        \footnotesize
         \begin{block}{}
 		  \begin{itemize}
 		    \item Colocamos todos os elementos que efetuam hash para o mesmo endereço em uma \textbf{lista ligada}.
 		  \end{itemize}

 		  \end{block} 
 
        \column{0.7\linewidth}
        \centering
        \includegraphics[bb=0 0 774 365,scale=0.3]{figs/encadeamento.png}        
\end{columns}
\end{frame}

\begin{frame}
\begin{columns}
        \column{0.3\linewidth}
        \footnotesize
         \begin{block}{}
 		  \begin{itemize}
 		    \item Endereço $j$ contém um ponteiro para início da lista.
 		    \item Se caso estiver vazio, $j$ conterá nullo.
 		  \end{itemize}

 		  \end{block} 
 
        \column{0.7\linewidth}
        \centering
        \includegraphics[bb=0 0 774 365,scale=0.3]{figs/encadeamento.png}        
\end{columns}
\end{frame}

\begin{frame}
\begin{columns}
        \column{0.35\linewidth}
        \footnotesize
         \begin{block}{Inserção}
 		 \begin{description}
 		  \item[Pior Caso] $O(1)$ 
 		 \end{description}
 		  \end{block} 
        \begin{block}{Remoção}
 		 \begin{description}
 		  \item[Pior Caso] $O(1)$
 		 \end{description}
 		 \begin{itemize}
 		  \item Lista Duplamente Encadeada
 		  \item Método recebe como entrada um elemento $X$ e não sua chave $k$
 		  $$X.ant.prox = X.prox$$
 		  $$X.prox.ant = X.ant$$
 		 \end{itemize}

 		  \end{block} 
        
        \column{0.7\linewidth}
        \centering
        \includegraphics[bb=0 0 774 365,scale=0.3]{figs/encadeamento.png}        
\end{columns}
\end{frame}

\begin{frame}
\begin{columns}
        \column{0.35\linewidth}
        \footnotesize         
        \begin{block}{Pesquisa}
 		 \begin{description}
 		  \item[Pior Caso] $\Theta(n)$ 
 		 \end{description}
 		 \begin{itemize}
 		  \item As $\mathbf{n}$ chaves no mesmo endereço, criando uma lista de dimensão $\mathbf{d}$
 		 \end{itemize}

 		  \end{block} 
        \column{0.7\linewidth}
        \centering
        \includegraphics[bb=0 0 774 365,scale=0.3]{figs/encadeamento.png}        
\end{columns}
\end{frame}

\begin{frame}
\begin{columns}
        \column{0.35\linewidth}
        \footnotesize         
        \begin{block}{Pesquisa}
 		 \begin{description}
 		  \item[Sucesso] $\Theta(1+\alpha)$
 		  \item[Sem Sucesso] $\Theta(1+\alpha)$
 		 \end{description}
 		 \begin{itemize}
 		  \item Hash uniforme
 		 \end{itemize}

 		  \end{block} 
        \column{0.7\linewidth}
        \centering
        \includegraphics[bb=0 0 774 365,scale=0.3]{figs/encadeamento.png}        
\end{columns}
\end{frame}

\begin{frame}{Exercício C - Encadeamento}
\begin{itemize}
 \item Demostre a inserção das chaves \textbf{5,28,19,15,20,33,12,17,10} em uma tabela hash com colisões por encadeamento. Seja tabela com 9 posições e com função $h(k) = k \ \mathbf{mod} \ 9$  
\end{itemize}
\end{frame}

\begin{frame}{Exercício C - Encadeamento}
\begin{columns}
        \column{0.5\linewidth}
        \footnotesize         
        \begin{block}{}
 		 \begin{itemize}
 		  \item \textbf{5,28,19,15,20,33,12,17,10}
 		  \item $h(k) = k \ \mathbf{mod} \ 9$  
 		 \end{itemize}

 		  \end{block} 
        \column{0.5\linewidth}
    
        \begin{table}[]
        \centering
        \footnotesize
        \begin{tabular}{|l|l|}
        \hline
        h(k) & k \\ \hline
            &   \\ \hline
            &   \\ \hline
            &   \\ \hline
            &   \\ \hline
            &   \\ \hline
            &   \\ \hline
            &   \\ \hline
            &   \\ \hline
            &   \\ \hline
        \end{tabular}%

        \end{table}       
        \end{columns}
\end{frame}

\begin{frame}{Exercício C - Encadeamento}
\begin{columns}
        \column{0.5\linewidth}
        \footnotesize         
        \begin{block}{}
 		 \begin{itemize}
 		  \item \textbf{28,19,15,20,33,12,17,10}
 		  \item $h(k) = 5 \ \mathbf{mod} \ 9$  
 		 \end{itemize}

 		  \end{block} 
        \column{0.5\linewidth}
    
        \begin{table}[]
        \centering
        \footnotesize
        \begin{tabular}{|l|l|}
        \hline
        h(k) & k \\ \hline
            &   \\ \hline
            &   \\ \hline
            &   \\ \hline
            &   \\ \hline
            &   \\ \hline
        5   & 5   \\ \hline
            &   \\ \hline
            &   \\ \hline
            &   \\ \hline
        \end{tabular}%

        \end{table}       
        \end{columns}
\end{frame}



\begin{frame}{Exercício C - Encadeamento}
\begin{columns}
        \column{0.5\linewidth}
        \footnotesize         
        \begin{block}{}
 		 \begin{itemize}
 		  \item \textbf{19,15,20,33,12,17,10}
 		  \item $h(k) = 28 \ \mathbf{mod} \ 9$  
 		 \end{itemize}

 		  \end{block} 
        \column{0.5\linewidth}
    
        \begin{table}[]
        \centering
        \footnotesize
        \begin{tabular}{|l|l|}
        \hline
        h(k) & k \\ \hline
            &   \\ \hline
        1   & 28  \\ \hline
            &   \\ \hline
            &   \\ \hline
            &   \\ \hline
        5   & 5   \\ \hline
            &   \\ \hline
            &   \\ \hline
            &   \\ \hline
        \end{tabular}%

        \end{table}       
        \end{columns}
\end{frame}

\begin{frame}{Exercício C - Encadeamento}
\begin{columns}
        \column{0.5\linewidth}
        \footnotesize         
        \begin{block}{}
 		 \begin{itemize}
 		  \item \textbf{15,20,33,12,17,10}
 		  \item $h(k) = 19 \ \mathbf{mod} \ 9$  
 		 \end{itemize}

 		  \end{block} 
        \column{0.5\linewidth}
    
        \begin{table}[]
        \centering
        \footnotesize
        \begin{tabular}{|l|l|}
        \hline
        h(k) & k \\ \hline
            &   \\ \hline
        1   & 19$\rightarrow$28  \\ \hline
            &   \\ \hline
            &   \\ \hline
            &   \\ \hline
        5   & 5   \\ \hline
            &   \\ \hline
            &   \\ \hline
            &   \\ \hline
        \end{tabular}%

        \end{table}       
        \end{columns}
\end{frame}


\begin{frame}{Exercício C - Encadeamento}
\begin{columns}
        \column{0.5\linewidth}
        \footnotesize         
        \begin{block}{}
 		 \begin{itemize}
 		  \item \textbf{20,33,12,17,10}
 		  \item $h(k) = 15 \ \mathbf{mod} \ 9$  
 		 \end{itemize}

 		  \end{block} 
        \column{0.5\linewidth}
    
        \begin{table}[]
        \centering
        \footnotesize
        \begin{tabular}{|l|l|}
        \hline
        h(k) & k \\ \hline
            &   \\ \hline
        1   & 19$\rightarrow$28  \\ \hline
            &   \\ \hline
            &   \\ \hline
            &   \\ \hline
        5   & 5   \\ \hline
        6   & 15   \\ \hline
            &   \\ \hline
            &   \\ \hline
        \end{tabular}%

        \end{table}       
        \end{columns}
\end{frame}


\begin{frame}{Exercício C - Encadeamento}
\begin{columns}
        \column{0.5\linewidth}
        \footnotesize         
        \begin{block}{}
 		 \begin{itemize}
 		  \item \textbf{33,12,17,10}
 		  \item $h(k) = 20 \ \mathbf{mod} \ 9$  
 		 \end{itemize}

 		  \end{block} 
        \column{0.5\linewidth}
    
        \begin{table}[]
        \centering
        \footnotesize
        \begin{tabular}{|l|l|}
        \hline
        h(k) & k \\ \hline
            &   \\ \hline
        1   & 19$\rightarrow$28  \\ \hline
        2   & 20  \\ \hline
            &    \\ \hline
            &   \\ \hline
        5   & 5   \\ \hline
        6   & 15   \\ \hline
            &   \\ \hline
            &   \\ \hline
        \end{tabular}%

        \end{table}       
        \end{columns}
\end{frame}


\begin{frame}{Exercício C - Encadeamento}
\begin{columns}
        \column{0.5\linewidth}
        \footnotesize         
        \begin{block}{}
 		 \begin{itemize}
 		  \item \textbf{12,17,10}
 		  \item $h(k) = 33 \ \mathbf{mod} \ 9$  
 		 \end{itemize}

 		  \end{block} 
        \column{0.5\linewidth}
    
        \begin{table}[]
        \centering
        \footnotesize
        \begin{tabular}{|l|l|}
        \hline
        h(k) & k \\ \hline
            &   \\ \hline
        1   & 19$\rightarrow$28  \\ \hline
        2   & 20  \\ \hline
            &    \\ \hline
            &   \\ \hline
        5   & 5   \\ \hline
        6   & 33 $\rightarrow$ 15   \\ \hline
            &   \\ \hline
            &   \\ \hline
        \end{tabular}%

        \end{table}       
        \end{columns}
\end{frame}


\begin{frame}{Exercício C - Encadeamento}
\begin{columns}
        \column{0.5\linewidth}
        \footnotesize         
        \begin{block}{}
 		 \begin{itemize}
 		  \item \textbf{17,10}
 		  \item $h(k) = 12 \ \mathbf{mod} \ 9$  
 		 \end{itemize}

 		  \end{block} 
        \column{0.5\linewidth}
    
        \begin{table}[]
        \centering
        \footnotesize
        \begin{tabular}{|l|l|}
        \hline
        h(k) & k \\ \hline
            &   \\ \hline
        1   & 19$\rightarrow$28  \\ \hline
        2   & 20  \\ \hline
        3   & 12   \\ \hline
            &    \\ \hline
        5   & 5   \\ \hline
        6   & 33 $\rightarrow$ 15   \\ \hline
            &   \\ \hline
            &   \\ \hline
        \end{tabular}%

        \end{table}       
        \end{columns}
\end{frame}


\begin{frame}{Exercício C - Encadeamento}
\begin{columns}
        \column{0.5\linewidth}
        \footnotesize         
        \begin{block}{}
 		 \begin{itemize}
 		  \item \textbf{10}
 		  \item $h(k) = 17 \ \mathbf{mod} \ 9$  
 		 \end{itemize}

 		  \end{block} 
        \column{0.5\linewidth}
    
        \begin{table}[]
        \centering
        \footnotesize
        \begin{tabular}{|l|l|}
        \hline
        h(k) & k \\ \hline
            &   \\ \hline
        1   & 19$\rightarrow$28  \\ \hline
        2   & 20  \\ \hline
        3   & 12   \\ \hline
            &    \\ \hline
        5   & 5   \\ \hline
        6   & 33 $\rightarrow$ 15   \\ \hline
            &   \\ \hline
        8   & 17  \\ \hline
        \end{tabular}%

        \end{table}       
        \end{columns}
\end{frame}

\begin{frame}{Exercício C - Encadeamento}
\begin{columns}
        \column{0.5\linewidth}
        \footnotesize         
        \begin{block}{}
 		 \begin{itemize}
 		  \item \textbf{$\oslash$}
 		  \item $h(k) = 10 \ \mathbf{mod} \ 9$  
 		 \end{itemize}

 		  \end{block} 
        \column{0.5\linewidth}
    
        \begin{table}[]
        \centering
        \footnotesize
        \begin{tabular}{|l|l|}
        \hline
        h(k) & k \\ \hline
            &   \\ \hline
        1   & 10 $\rightarrow$ 19$\rightarrow$28  \\ \hline
        2   & 20  \\ \hline
        3   & 12   \\ \hline
            &    \\ \hline
        5   & 5   \\ \hline
        6   & 33 $\rightarrow$ 15   \\ \hline
            &   \\ \hline
        8   & 17  \\ \hline
        \end{tabular}%

        \end{table}       
        \end{columns}
\end{frame}


\begin{frame}{Exercício C - Aluno - Encadeamento}
\begin{itemize}
 \item Forneça o conteúdo da \textit{hash table} resultante quando você insere itens com as chaves ATLAS nessa ordem em uma tabela inicialmente vazia de M = 5, usando encadeamento com listas não ordenadas. Use a função hash $11k \ \mathbf{mod}\ M$ para transformar a k-ésima letra do alfabeto em um índice de tabela.
 \item $A = 1, B = 2, ... $
\end{itemize}

\end{frame}

\begin{frame}[fragile]{Implementação em C - Encadeamento}
\footnotesize
\begin{lstlisting}[style=CStyle]
#include <stdio.h>
#include <stdlib.h>

#define SIZE 10

// Estrutura para o Nodo ou Elemento da lista encadeada
typedef struct Element {
    int key;
    struct Element* next;
} Element;

// Estrutura para a tabela de hash
typedef struct HashTable {
    Element** table;//Ponteiro para Ponteiro
} HashTable;

\end{lstlisting}
\end{frame}

\begin{frame}[fragile]{Implementação em C - Encadeamento}
\footnotesize
\begin{lstlisting}[style=CStyle]
// Funcao para criar um novo elemento
Element* createElement(int key) {
    Element* newElement = (Element*)malloc(sizeof(Element));
    newElement->key = key;
    newElement->next = NULL;
    return newElement;
}
// Funcao para criar a tabela de hash
HashTable* createHashTable() {
    HashTable* hashTable = (HashTable*)malloc(sizeof(HashTable));
    hashTable->table =(Element**)malloc(SIZE *sizeof(Element*));
    // Inicializar cada posicao da tabela com NULL
    for (int i = 0; i < SIZE; i++) {
        hashTable->table[i] = NULL;
    }
    return hashTable;
}

\end{lstlisting}
\end{frame}

\begin{frame}[fragile]{Implementação em C - Encadeamento}
\footnotesize
\begin{lstlisting}[style=CStyle]
// Funcao hash
int hashFunction(int key) {
    return key % SIZE;
}
\end{lstlisting}
\end{frame}


\begin{frame}[fragile]{Implementação em C - Encadeamento}
\footnotesize
\begin{lstlisting}[style=CStyle]
void insert(HashTable* hashTable, int key) {
    int index = hashFunction(key);
    // Criar um novo elemento ou nodo
    Element* newElement = createElement(key);
    // Inserir o elemento na posicao correspondente da tabela
    if (hashTable->table[index] == NULL) {
        // Caso a posicao esteja vazia, o novo elemento se torna a cabeca da lista
        hashTable->table[index] = newElement;
    } else {//Senao, add o novo elemento ao final da lista
        Element* currentElement = hashTable->table[index];
        while (currentElement->next != NULL) {
            currentElement = currentElement->next;
        }
        currentElement->next = newElement;
    }
}
\end{lstlisting}
\end{frame}

\begin{frame}[fragile]{Implementação em C - Encadeamento}
\footnotesize
\begin{lstlisting}[style=CStyle]
// Funcao para imprimir a tabela de hash
void printHashTable(HashTable* hashTable) {
    for (int i = 0; i < SIZE; i++) {
        printf("Posicao %d: ", i);
        Element* currentElement = hashTable->table[i];
        while (currentElement != NULL) {
            printf("%d ", currentElement->key);
            currentElement = currentElement->next;
        }
        printf("\n");
    }
}
\end{lstlisting}
\end{frame}

\begin{frame}[fragile]{Implementação em C - Encadeamento}
\footnotesize
\begin{lstlisting}[style=CStyle]
int main() {
    // Criar a tabela de hash
    HashTable* hashTable = createHashTable();

    // Inserir elementos na tabela de hash
    insert(hashTable, 7);
    insert(hashTable, 12);
    insert(hashTable, 23);
    insert(hashTable, 35);
    insert(hashTable, 14);
    insert(hashTable, 13);
    insert(hashTable, 33);
    insert(hashTable, 40);
    insert(hashTable, 12);

    // Imprimir a tabela de hash
    printHashTable(hashTable);

    return 0;
}
\end{lstlisting}
\end{frame}

\begin{frame}[fragile]{Implementação em C - Encadeamento}
\begin{center}
 \includegraphics[bb=0 0 481 316,scale=0.5]{figs/printHashingChainInC.png}
 % printHashingChainInC.png: 481x316 px, 72dpi, 16.97x11.15 cm, bb=0 0 481 316
\end{center}

\end{frame}


\begin{frame}[fragile]{Implementação em Python} 
\centering\footnotesize
\begin{lstlisting}[style=PythonStyle]
class HashingChain:
    

    # Funcao add valor na Tabela hash
    def insere(self,hashtable, valor_chave, valor):
        if(self.funcao == "Divisao"):
            endereco = self.hashingDivisao(valor_chave,hashtable)
        elif (self.funcao == "Multiplicacao"):
            endereco = self.hashingMultiplicacao(valor_chave,hashtable)
        hashtable[endereco].insert(0,valor) 
\end{lstlisting}


\end{frame} 

\begin{frame}[fragile]{Implementação em Python} 
\centering\footnotesize
\begin{lstlisting}[style=PythonStyle]
class HashingChain:
    def display_hash(self,hashTable):
        for i in range(len(hashTable)):
            print(i, end=" ")
            for j in hashTable[i]:
                print("-->", end=" ")
                print(j, end=" ")
            print()
\end{lstlisting}


\end{frame}

\begin{frame}[fragile]{Implementação em Python} 
\footnotesize
\begin{lstlisting}[style=PythonStyle]
if __name__ == '__main__':
    HashTable = [[] for x in range(9)]
    hashing = HashingChain(funcao="Divisao")
    for valor in [5,28,19,15,20,33,12,17,10]:
        hashing.insere(HashTable, valor, str(valor))
    hashing.display_hash(HashTable)
\end{lstlisting}
\begin{block}{Resultado}
$0$\\ 
$1 --> 10 --> 19 --> 28$\\ 
$2 --> 20$\\ 
$3 --> 12$\\ 
$4$\\ 
$5 --> 5$\\ 
$6 --> 33 --> 15$\\ 
$7$\\ 
$8 --> 17$ 
\end{block}
\end{frame}

% \subsection{Trabalho III - Encadeamento}
% \begin{frame}
% \begin{itemize}
%  \item Existe a hipótese de que é possível obter ganhos substanciais de desempenho se modificarmos o esquema de encadeamento de tal modo que cada lista seja mantida em sequencia ordenada. Como modificação afeta o tempo de execução para pesquisas bem sucedidas e sem sucesso, inserções e eliminações ?
%  \begin{itemize}
%   \item Implemente essa tabela hash:
%
%   \begin{itemize}
%   \item Implemente lista ordenada
%   \item chaves são alphanuméricas
%   \item Verifique em termos de notação assintótica considerando tamanho da lista.
%   \end{itemize}
%
%  \end{itemize}
%
% \end{itemize}
%
% \end{frame}

\subsection{Endereçamento Aberto}
\begin{frame}
 \begin{itemize}
  \item Os elementos estão armazenados na própria tabela hash.
  \item Não existem listas ou estruturas externas.
  \item Tabela pode ficar ``cheia''  
  \item $\alpha \leq 1$
  \item Não usa ponteiros
  \begin{itemize}
   \item Calculamos a sequência de endereços
   \item Memória extra liberada
   \item Menor número de colisões
   \item Recuperação mais rápida
  \end{itemize}
 \end{itemize}
\end{frame}

\begin{frame}
 \begin{itemize}
  \item \emph{Possível fornecer vários endereços a partir de uma única chave}
  \item Processa uma análise sequêncial ou sondagem nos endereços
  \begin{itemize}
   \item $h: U \times \{0,1,..,m-1\} \rightarrow \{0,1,..,m-1\}$
   \item $h(k,x)$
  \end{itemize}
  \item Pesquisar uma chave $k$
  \begin{itemize}
   \item $h(k,0)\rightarrow h(k,1)\rightarrow ... \rightarrow h(k,m-1)$
   \begin{itemize}
    \item Permutação de endereços para cada chave $\langle h(k,0), h(k,1),...,h(k,m-1)\rangle$
   \end{itemize}
   \item $x=0,...,m-1$
   \begin{itemize}
    \item sequência de tentativas ou sondagem
   \end{itemize}
   
  \end{itemize}
\item sequência uniforme
\begin{itemize}
 \item qualquer $\mathbf{m!}$ dos endereços, têm igual probabilidade de ser produzida por $\mathbf{h}$
\end{itemize}

  \end{itemize}

\end{frame}



\begin{frame}
 \begin{itemize}
  \item Operação de remoção 
  \begin{itemize}
   \item Romperia a sequência de tentativas
   \begin{itemize}
    \item Tornaria impossível recuperar qualquer chave $\mathbf{k}$
    \item Encontraria \textbf{null} no endereço
   \end{itemize}

   \item Assinalar a posição do armazenamento: 
   \begin{description}
    \item[Livre (L)] Quando posição não foi utilizada
    \item[Ocupado (O)] Quando uma chave está armazenada
    \item[Removido (R)] Quando armazena uma chave que já foi removida - uma nova chave poderá ocupar a posição marcada como removida
   \end{description}
  \end{itemize}
  \end{itemize}

\end{frame}


\begin{frame}
 \begin{itemize}
  \item Existem três métodos para a determinação da sequência de tentativas ou sondagem $h(x,k)\ \forall \ k = 0,1,...,m-1$.
  \begin{itemize}
   \item tentativa linear 
   \item tentativa quadrática
   \item hash duplo
  \end{itemize}

 \end{itemize}

\end{frame}


\begin{frame}{Tentativa linear}
\begin{itemize}
 \item Dada função hash auxiliar $h':U\rightarrow {0,1,...,m-1}$, método de tentativa linear usa a função hash:
 $$h(k,i) = (h'(k) + i) \ mod \ m$$
 \item $i=0,1,...,m-1$
 \item A ideia consiste em tentar armazenar chave $k$ em $h'(k)$, se este já está ocupado, tentar no endereço consecutivo $h'(k) + 1$, até encontrar posição vazia.
\item Problema: \textbf{Agrupamento primário}
\begin{itemize}                                     \item Longas sequência de posições ocupadas são construídas, aumentando o tempo médio de pesquisa. 
\item Causa: Uma posição vazia precedida por \textbf{i} posições completas é preenchida em seguida com probabilidade $\dfrac{i+1}{m}$
\end{itemize}
\end{itemize}
\end{frame}



\begin{frame}{Tentativa linear}
\begin{columns}
        \column{0.8\linewidth}
        \footnotesize         
        \begin{block}{}
 		 \begin{itemize}
 		  \item $m = 11$ e $h'(x) = x \ mod \ 11 $	
 		  \item $h(k,i) = (h'(k) + i) \ mod \ m$
 		  \item $[20,30,2,13,25,24,10,9]$
 		  \begin{itemize}
 		   \item $h(20,0) = h'(20)\ mod\ 11 = \mathbf{9}$
 		   \item $h(30,0) = h'(30)\ mod\ 11 = \mathbf{8}$
 		   \item $h(2,0) =  h'(2)\  mod\ 11 = \mathbf{2}$
 		   \item $h(13,0) =  h'(13)\  mod\ 11 = 2$ 
 		   \begin{itemize}
 		    \item $h(13,1) =  (h'(13)+1)\  mod\ 11 = \mathbf{3}$ 
 		   \end{itemize}

 		   \item $25\ mod\ 11 = 3\rightarrow 3+ 1 = \mathbf{4}$
 		   \item $24\ mod\ 11 = 2 \rightarrow 2+1 ... \rightarrow 2+3=\mathbf{5}$
 		   \item $10\ mod\ 11 = \mathbf{10}$
 		   \item $9\ mod\ 11 = 9\rightarrow 9+1\rightarrow 9+2=0$
 		  \end{itemize}

 		 \end{itemize}

 		  \end{block} 
        \column{0.5\linewidth}    
        \centering
        \includegraphics[bb=0 0 121 375,scale=.5]{figs/tentativaLinear.png}             
        \end{columns}
        
\end{frame}

\begin{frame}[fragile]{Implementação em C}
%\centering\footnotesize
\begin{lstlisting}[style=CStyle]

int tentativaLinear(int table[],int hash)
{
    // Tratamento de colisao por sondagem linear
    while (table[hash] != -1) {
        hash = (hash + 1) % TABLE_SIZE;
    }
    return hash;
}
\end{lstlisting}
\end{frame}

\begin{frame}[fragile]{Implementação em C}
\footnotesize
\begin{lstlisting}[style=CStyle]
void insere(int table[], int key, char funcao) {
    if(funcao == 'D'){
        int hash = key % TABLE_SIZE;
        hash = tentativaLinear(table,hash)
        table[hash] = key;
    }
    else if(funcao == 'M')
    {
        float A = (sqrt(5)-1)/2;
        float m = TABLE_SIZE;
        int hash = floor(m*(fmod(key*A , 1)));
        table[hash] = key;
    }
}
\end{lstlisting}
\end{frame}


\begin{frame}[fragile]{Implementação em C}
\footnotesize
\begin{lstlisting}[style=CStyle]
#define TABLE_SIZE 11

int main() {
    int hash_table[TABLE_SIZE];
    initialize_table(hash_table, TABLE_SIZE);
    int keys[] = {20,30,2,13,25,24,10,9};
    int num_keys = sizeof(keys) / sizeof(keys[0]);
    for (int i = 0; i < num_keys; i++) {
        int key = keys[i];
        insere(hash_table, key,'D');
    }
    // Imprimir tabela hash
    for (int i = 0; i < TABLE_SIZE; i++) {
        printf("Index: %d, Key: %d\n", i, hash_table[i]);
    }
    return 0;
}
\end{lstlisting}
\end{frame}

\begin{frame}[fragile]{Implementação em Python} 
\centering\footnotesize
\begin{lstlisting}[style=PythonStyle]
  class HashingOpenAddressing:
    def __init__(self,metodo):
        self.metodo = metodo;

    def linear(self, chave, hashTable):
    #HashTable = np.array([[None] * 11, ['L'] * 11])
        m = len(hashTable[0])
        h1 = chave % m
        for i in range(m):
            h = (h1+i) %  m
            if(hashTable[1,h] == 'L'):
                break;
        return h;
\end{lstlisting}
\end{frame}

\begin{frame}[fragile]{Implementação em Python} 
\centering\footnotesize
\begin{lstlisting}[style=PythonStyle]
  # Funcao add valor na Tabela hash
    def insere(self,hashtable, valor_chave, valor):
        if(self.metodo == "linear"):
            endereco = self.linear(valor_chave,hashtable)
        elif(self.metodo == "quadratica"):
            endereco = self.quadratica(valor_chave, hashtable)
        hashtable[0, endereco] = valor
        hashtable[1, endereco] = 'O'
    
    def display_hash(self, hashTable):
        for i in range(len(hashTable[0])):
            print(str(i) + " --> " + str(hashTable[0,i])+" ("+str(hashTable[1,i])+")", end="\n")
\end{lstlisting}
\end{frame}

\begin{frame}[fragile]{implementação em Python} 
\centering\footnotesize
\begin{lstlisting}[style=PythonStyle]
    HashTable = np.array([[None] * 11, ['L'] * 11])
    hashing = HashingOpenAddressing(metodo="linear")
    for valor in [20, 30, 2, 13, 25, 24, 10, 9]:
        hashing.insere(HashTable, valor, str(valor))
    hashing.display_hash(HashTable)
\end{lstlisting}
\begin{block}{Resultado}
\tiny
$0 --> 9 (O)$\\
$1 --> None (L)$\\
$2 --> 2 (O)$\\
$3 --> 13 (O)$\\
$4 --> 25 (O)$\\
$5 --> 24 (O)$\\
$6 --> None (L)$\\
$7 --> None (L)$\\
$8 --> 30 (O)$\\
$9 --> 20 (O)$\\
$10 --> 10 (O)$\\
\end{block}

\end{frame}

\begin{frame}{Tentativa Quadrática}
\begin{itemize}
 \item \textbf{Tentativa linear} as chaves tendem a se concentrar, criando \textbf{agrupamentos primários}, que aumentam muito o tempo de busca
 \begin{itemize}
  \item A ideia é obter sequências de endereços diversas para endereços próximos.
 \end{itemize}
 $$h(k,i) = (h'(k) + c_1i + c_2i^2)\ mod \ m$$
 \item $c1$ e $c2 \neq 0$
 \item $i  = 0,1,...,m-1$

\end{itemize}
\end{frame}

\begin{frame}{Tentativa Quadrática}
\begin{itemize}
 \item Esse método consegue evitar os agrupamentos primários da tentativa linear.
 \item \textbf{Agrupamento Secundário}
 \begin{itemize}
  \item $h(k_1,0) = h(k_2,0) \rightarrow h(k_1,i) = h(k_2,i)$
  \item Degradações são reduzidas se comparadas com a Tentaiva linear.
 \end{itemize}
\end{itemize}
\end{frame}

\begin{frame}{Tentativa Quadrática}
\begin{itemize}
 \item Os valores $m,c_1$ e $c_2$ devem ser escolhidos de tal forma que os endereços $h(x, k)$ correspondam a varrer toda a tabela para $k=0,1,...,m-1$
 \item As equações recorrentes fornecem uma maneira de calcular, diretamente esses endereços:
 $$h(x,0) = h'(x)$$
 $$h(x,k) = (h(x,k-1) + k) \ mod \ m, 0 < k < m$$
 \end{itemize}
\end{frame}

\begin{frame}{Tentativa Quadrática}
\begin{figure}[h]
 \centering
 \includegraphics[bb=0 0 172 100,scale=1.2]{figs/linear_quadratico.png}
 % linear_quadratico.png: 718x415 px, 300dpi, 6.08x3.51 cm, bb=0 0 172 100
 \caption*{Comparação de Comportamento}
\end{figure}

\end{frame}

\begin{frame}{Tentativa Quadrática}
\begin{columns}
        \column{0.8\linewidth}
        \footnotesize         
        \begin{block}{}
 		 \begin{itemize}
 		  \item $m = 11$ e $h'(x) = x \ mod \ 11 $ 		 
 		  \item $[20,30,2,13,25,24,10,9]$
 		  \begin{itemize}
 		   \item $h(20,0) = h'(20)\ mod\ 11 = \mathbf{9}$
 		   \item $h(30,0) = h'(30)\ mod\ 11 = \mathbf{8}$
 		   \item $h(2,0) =  h'(2)\  mod\ 11 = \mathbf{2}$
 		   \item $h(13,0) =  h'(13)\  mod\ 11 = 2$ 
 		   \begin{itemize}
 		    \item $h(13,1) =  (h'(13)+1^2)\  mod\ 11 = \mathbf{3}$ 
 		   \end{itemize}
 		   \item $25\ mod\ 11 = 3\rightarrow 3+ 1 = \mathbf{4}$
 		   \item $24\ mod\ 11 = 2 \rightarrow 2+1^2 \rightarrow 2+2^2=\mathbf{6}$
 		   \item $10\ mod\ 11 = \mathbf{10}$
 		   \item $9\ mod\ 11 = 9\rightarrow 9+1^2\rightarrow 9+2^2\rightarrow 9+3^2=\mathbf{7}$
 		  \end{itemize}

 		 \end{itemize}

 		  \end{block} 
        \column{0.5\linewidth}    
        \centering
        \includegraphics[bb=0 0 36 112,scale=1.3]{figs/tentativaQuadratica.png}
% tentativaQuadratica.png: 149x465 px, 300dpi, 1.26x3.94 cm, bb=0 0 36 112
            
        \end{columns}
        
\end{frame}


\begin{frame}[fragile]{Implementação em C}
%\centering\footnotesize
\begin{lstlisting}[style=CStyle]

int tentativaQuadratica(int table[],int hash)
{
    // Tratamento de colisao por sondagem quadratica
    for(int i = 0 ;i<TABLE_SIZE; i++)
    {
        int h = (hash + i*i) % TABLE_SIZE;
        if(table[h] == -1)
            return h;
    }
}
\end{lstlisting}
\end{frame}

\begin{frame}[fragile]{Implementação em C}
%\footnotesize
\begin{lstlisting}[style=CStyle]
void insere(int table[], int key, char funcao) {
    if(funcao == 'D'){
        int hash = key % TABLE_SIZE;
        hash = tentativaQuadratica(table,hash);
        table[hash] = key;
    }
    else if(funcao == 'M')
    {
        float A = (sqrt(5)-1)/2;
        float m = TABLE_SIZE;
        int hash = floor(m*(fmod(key*A , 1)));
        table[hash] = key;
    }
}
\end{lstlisting}
\end{frame}

\begin{frame}[fragile]{Implementação em C}
\footnotesize
\begin{lstlisting}[style=CStyle]
#define TABLE_SIZE 11
int main() {
    int hash_table[TABLE_SIZE];
    initialize_table(hash_table, TABLE_SIZE);
    int keys[] = {20,30,2,13,25,24,10,9};
    int num_keys = sizeof(keys) / sizeof(keys[0]);

    for (int i = 0; i < num_keys; i++) {
        int key = keys[i];
        insere(hash_table, key,'D');
    }
    for (int i = 0; i < TABLE_SIZE; i++) {
        printf("Index: %d, Key: %d\n", i, hash_table[i]);
    }

    return 0;
}

\end{lstlisting}
\end{frame}


\begin{frame}[fragile]{Implementação em Python} 
\centering\footnotesize
\begin{lstlisting}[style=PythonStyle]
  class HashingOpenAddressing:
    def __init__(self,metodo):
        self.metodo = metodo;

    def quadratica(self, chave, hashTable):
        m = len(hashTable[0])
        h1 = chave % m
        for i in range(m):
            h = (h1+i*i) %  m
            if(hashTable[1,h] == 'L'):
                break;
        return h;
\end{lstlisting}
\end{frame}

\begin{frame}[fragile]{implementação em Python} 
\centering\footnotesize
\begin{lstlisting}[style=PythonStyle]
HashTable = np.array([[None] * 11, ['L'] * 11])
hashing = HashingOpenAddressing(metodo="quadratica")
for valor in [20, 30, 2, 13, 25, 24, 10, 9]:
    hashing.insere(HashTable, valor, str(valor))
hashing.display_hash(HashTable)
\end{lstlisting}
\begin{block}{Resultado}
\tiny
$0 --> None (L)$\\
$1 --> None (L)$\\
$2 --> 2 (O)$\\
$3 --> 13 (O)$\\
$4 --> 25 (O)$\\
$5 --> None (L)$\\
$6 --> 24 (O)$\\
$7 --> 9 (O)$\\
$8 --> 30 (O)$\\
$9 --> 20 (O)$\\
$10 --> 10 (O)$\\
\end{block}
\end{frame}


\begin{frame}{Hash Duplo}
\begin{itemize}
 \item Um dos melhores métodos para endereçamento aberto
 \item Por que as permutações produzidas possuem características de permutações aleatórias
$$h(k,i) = (h_1(k) + i\times h_2(k))\ mod \ m$$
\item Endereço inicial $T[h_1(k)]$
\item Endereços sucessivos dependem da quantidade $h_2(k)$
\item O deslocamento depende da chave $k$ de duas maneiras
 \end{itemize}
\end{frame}

\begin{frame}{Hash Duplo}
\begin{columns}
        \column{0.7\linewidth}
        \footnotesize         
        \begin{block}{}
 		 \begin{itemize}
 		  \item $m = 13$ 
 		  \item $h_1(k) = k \ mod  \ 13$ \\ $h_2(k) = 1+ (k \ mod  \ 11)$ 	
 		  \item $h(k,i) = ( h_1(k) + i \times h_2(k))\ mod \ 13 $
 		  \item Inserir \textbf{k=14}
 		 \end{itemize}
%(14 mod 13 + 2×(1+(14 mod 11))) mod 13
 		  \end{block} 
        \column{0.5\linewidth}    
        \centering
        \includegraphics[bb=0 0 113 322,scale=0.5]{figs/hashduplo.png}            
        \end{columns}

\end{frame}


\begin{frame}{Hash Duplo}
\begin{columns}
        \column{0.7\linewidth}
        \footnotesize
        \begin{block}{}
 		 \begin{itemize}
 		  \item $m = 13$
 		  \item $h_1(k) = k \ mod  \ 13$ \\ $h_2(k) = 1+ (k \ mod  \ 11)$
 		  \item $h(k,i) = ( h_1(k) + i \times h_2(k))\ mod \ 13 $
 		  \item Inserir \textbf{k=14}
 		  \item $((14\ mod\ 13)+0×(1+(14\ mod\ 11)))\ mod\ 13 = 1$
 		 \end{itemize}
%(14 mod 13 + 2×(1+(14 mod 11))) mod 13
 		  \end{block}
        \column{0.5\linewidth}
        \centering
        \includegraphics[bb=0 0 113 322,scale=0.5]{figs/hashduplo.png}
        \end{columns}

\end{frame}

\begin{frame}{Hash Duplo}
\begin{itemize}
 \item Análise do hash do endereço aberto
 \begin{itemize}
  \item \emph{Hash uniforme}
  \item Em termos de $\alpha \leq 1$ e $n \leq m$ 
 \end{itemize}
\item Pesquisa
\begin{itemize}
 \item  malsucedida é executada $\dfrac{1}{1-\alpha}$
\item bem sucedida é executada $\dfrac{1}{\alpha}\ln \dfrac{1}{1-\alpha}$
\end{itemize}
 \end{itemize}
\end{frame}

\section{Hashing na Prática: Tópicos Avançados}
\subsection{Hashing Criptográfico e Segurança}
\begin{frame}{Hash criptográfico vs. não criptográfico}
  \begin{itemize}
    \item As funções vistas até aqui (divisão, multiplicação) priorizam velocidade -- perfeitas para tabelas hash em memória, mas inadequadas quando um adversário escolhe as chaves.
    \item Um hash \textbf{criptográfico} (ex.: SHA-256) acrescenta três garantias que o hash de tabela não precisa oferecer:
    \begin{itemize}
      \item resistência à pré-imagem: dado $h(x)$, é inviável achar $x$;
      \item resistência à segunda pré-imagem: dado $x$, é inviável achar $y \ne x$ com $h(y)=h(x)$;
      \item resistência a colisões: é inviável achar qualquer par $x \ne y$ com $h(x)=h(y)$.
    \end{itemize}
    \item Hashes não criptográficos (MurmurHash, xxHash, FNV) são muito mais rápidos, mas colidir propositalmente é trivial -- e é exatamente aí que mora o risco.
  \end{itemize}
\end{frame}

\begin{frame}{Ataques HashDoS e SipHash}
  \begin{itemize}
    \item Se um serviço web usa uma função hash previsível para indexar parâmetros recebidos (formulários, JSON), um atacante pode submeter milhares de chaves que colidem propositalmente.
    \item Toda inserção degrada de $O(1)$ para $O(n)$: uma tabela hash inteira vira, na prática, uma lista encadeada -- o \textbf{HashDoS} (negação de serviço algorítmica), relatado publicamente em 2011 contra várias linguagens de script.
    \item Mitigação adotada por Python (\texttt{dict}, desde 3.3), Rust e Ruby: usar \textbf{SipHash}, uma função pseudoaleatória rápida, com uma \emph{seed} de 128 bits sorteada aleatoriamente a cada execução do processo.
    \begin{itemize}
      \item Sem conhecer a seed, o atacante não consegue prever colisões.
      \item É o mesmo princípio de projeto por trás do ASLR: aleatoriedade em tempo de execução neutraliza um ataque que dependeria de determinismo.
    \end{itemize}
  \end{itemize}
\end{frame}

\subsection{Cuckoo Hashing e Robin Hood Hashing}
\begin{frame}{Cuckoo hashing}
  \begin{columns}[T]
    \column{.55\textwidth}
    \begin{itemize}
      \item Usa duas tabelas (ou uma tabela com duas funções hash, $h_1$ e $h_2$).
      \item Ao inserir $x$: se $T_1[h_1(x)]$ está livre, ocupa; senão, \textbf{expulsa} o ocupante atual e o realoja em $T_2[h_2(\cdot)]$ -- podendo expulsar outro elemento em cascata, como um filhote de cuco no ninho alheio.
      \item Garante busca e remoção em $O(1)$ \textbf{no pior caso}: cada chave só pode estar em duas posições possíveis.
      \item Ciclos de expulsão longos demais disparam um \emph{rehash} completo com novas funções.
    \end{itemize}
    \column{.4\textwidth}\centering\scriptsize
    \begin{tikzpicture}[x=9mm,y=8mm]
      \node[slot] (a) at (0,1) {$y$};
      \node[slot] (b) at (0,0) {$x$};
      \node[slot] (c) at (2,.5) {$y$};
      \draw[redarrow] (a) -- node[above,font=\tiny]{expulso} (c);
      \node[below=2mm of b,align=center] {$x$ ocupa\\ o lugar de $y$};
    \end{tikzpicture}
  \end{columns}
\end{frame}

\begin{frame}{Robin Hood hashing}
  \begin{itemize}
    \item Variante de endereçamento aberto (como a tentativa linear), mas com uma regra de troca justa.
    \item Cada elemento guarda sua \textbf{distância até o endereço ideal} $h'(k)$ (quantas sondagens foram necessárias).
    \item Ao inserir, se o elemento sendo colocado já percorreu mais distância que o ocupante atual da posição, eles trocam de lugar -- ``tira-se do rico (pouca distância) para dar ao pobre (muita distância)''.
    \item Resultado: a variância das distâncias de sondagem cai bastante, tornando o pior caso muito mais previsível que na tentativa linear pura.
    \item É a técnica por trás do \texttt{HashMap} padrão do Rust (implementação \emph{hashbrown}, inspirada nas SwissTables do Google), combinada com SipHash para resistir a HashDoS.
  \end{itemize}
\end{frame}

\subsection{Consistent Hashing}
\begin{frame}{Consistent hashing: motivação}
  \begin{itemize}
    \item No hashing tradicional, $h(k) = k \bmod m$: mudar $m$ (adicionar ou remover um servidor) remapeia \emph{quase todas} as chaves.
    \item Em um sistema distribuído com $m$ servidores de cache ou de armazenamento, isso significa invalidar quase tudo a cada mudança de cluster -- inviável em produção.
    \item \textbf{Consistent hashing} mapeia servidores \emph{e} chaves no mesmo espaço circular $[0, 2^{k}-1]$ (o ``anel'').
    \item Cada chave pertence ao primeiro servidor encontrado ao percorrer o anel no sentido horário.
  \end{itemize}
  \centering\small
  \begin{tikzpicture}[scale=.62]
    \draw[thick] (0,0) circle (1.3);
    \foreach \a/\lab in {30/{$S_1$},150/{$S_2$},260/{$S_3$}}{
      \node[slot,minimum width=6mm,minimum height=5mm,fill=hblockbg] at (\a:1.3) {\lab};
    }
    \foreach \a/\lab in {70/{$k_1$},190/{$k_2$},320/{$k_3$}}{
      \node[font=\tiny] at (\a:1.3) {$\bullet$};
      \node[font=\tiny] at (\a-8:1.65) {\lab};
    }
  \end{tikzpicture}
\end{frame}

\begin{frame}{Nós virtuais e uso em produção}
  \begin{itemize}
    \item Remover ou adicionar um servidor remapeia só as chaves entre ele e o próximo nó do anel -- em média, $\frac{1}{m}$ das chaves, não todas.
    \item Problema prático: poucos servidores reais distribuem o anel de forma desigual. Solução: cada servidor físico recebe várias posições virtuais (\emph{virtual nodes}) espalhadas no anel, equilibrando a carga.
    \item Usado, entre outros, por:
    \begin{itemize}
      \item Amazon DynamoDB e Apache Cassandra, para particionar dados entre nós;
      \item redes de CDN e Memcached distribuído, para decidir qual servidor de cache atende cada chave.
    \end{itemize}
  \end{itemize}
\end{frame}

\subsection{Bloom Filters}
\begin{frame}{Bloom filters}
  \small
  \begin{columns}[T]
    \column{.57\textwidth}
    \begin{itemize}
      \item Estrutura probabilística para testar pertencimento a um conjunto, com muito menos memória que guardar os elementos.
      \item Um vetor de $m$ bits (zerado) e $k$ funções hash independentes.
      \item \textbf{Inserir} $x$: liga os bits $h_1(x),\ldots,h_k(x)$.
      \item \textbf{Consultar} $x$: algum bit em 0 $\Rightarrow$ $x$ certamente \emph{não} está; todos em 1 $\Rightarrow$ $x$ \emph{provavelmente} está.
    \end{itemize}
    \column{.4\textwidth}\centering\scriptsize
    \begin{tikzpicture}[x=5.6mm,y=5.6mm]
      \foreach \v/\i in {0/0,1/1,0/2,1/3,1/4,0/5,1/6,0/7}{
        \ifnum\v=1 \node[cellG] at (\i,0) {1}; \else \node[cellW] at (\i,0) {0}; \fi
      }
    \end{tikzpicture}
  \end{columns}
  \medskip
  \begin{itemize}
    \item Nunca há falso negativo, mas pode haver \textbf{falso positivo} quando bits de elementos diferentes se sobrepõem.
    \item Taxa de falso positivo aproximada: $\left(1-e^{-kn/m}\right)^{k}$, para $n$ elementos inseridos.
  \end{itemize}
\end{frame}

\begin{frame}{Bloom filters: aplicações}
  \begin{itemize}
    \item \textbf{Bancos de dados} (Cassandra, HBase, RocksDB): antes de ler um arquivo em disco, consultam um Bloom filter em memória para saber se a chave \emph{pode} estar ali -- evitam leituras caras que dariam em vazio.
    \item \textbf{Navegadores}: listas de URLs maliciosas (Google Safe Browsing) usam Bloom filters para checar localmente, sem baixar a lista inteira.
    \item \textbf{Sistemas distribuídos e CDNs}: deduplicação e checagem rápida de cache antes de uma requisição de rede.
    \item O trade-off é sempre o mesmo: trocar uma pequena taxa de falsos positivos por uma economia enorme de memória e tempo.
  \end{itemize}
\end{frame}

% =============================================================================
\section{Conclusões}
\begin{frame}{Conclusões}
  \begin{itemize}
    \item Tabelas hash trocam a ordenação por um mapeamento direto chave $\to$ endereço, buscando $\Theta(1)$ em média.
    \item Colisões são inevitáveis; encadeamento e endereçamento aberto (linear, quadrático, hash duplo) resolvem-nas com trade-offs distintos de memória e agrupamento.
    \item O fator de carga $\alpha$ governa o desempenho em qualquer uma dessas estratégias.
    \item Técnicas modernas -- cuckoo/Robin Hood hashing, consistent hashing, Bloom filters, SipHash -- adaptam a mesma ideia central a garantias de pior caso, escala distribuída e segurança.
  \end{itemize}
\end{frame}


% \begin{frame}[fragile]{One-Way ANOVA  - Exemplo Prático}
%          \centering\footnotesize\begin{lstlisting}[language=R]
% Y <- c(rnorm(10,21,1),rnorm(10,17,2),rnorm(10,15,2)) 
% X <- rep(c("T1","T2","T3"),each=10) 
% dados <- data.frame(Y,X)
% summary(lm(Y ~ X))
%             \end{lstlisting}
%      \begin{itemize}
%          \item Resultado: 
%         \begin{figure}[!h]
%             \centering
%             \includegraphics[scale=0.25]{figs/logo.eps}
%             \label{fig:anovaResult}
%         \end{figure}
%      \end{itemize}
% 
% \end{frame} 


\section{Referências Bibliográficas}
\begin{frame}[t,allowframebreaks]
\frametitle{Referências}
%\bibliographystyle{amsalpha}
%\bibliography{bibfile}
\begin{thebibliography}{*}

        \bibitem[A01]{A01} Lee K.D., Hubbard S. (2015) Trees. In: Data Structures and Algorithms with Python. Undergraduate Topics in Computer Science. Springer, Cham. Retrieved from \url{https://doi.org/10.1007/978-3-319-13072-9_6}
        
        \bibitem[A02]{A02}Hubbard, J. (2007). Schaum’s Outline sof Data Structures with Java. Retrieved from \url{http://www.amazon.com/Schaums-Outline-Data-Structures-Java/dp/0071476989}
        
        \bibitem[A03]{A03} Cormen, T. H., Leiserson, C. E., \& Stein, R. L. R. E. C. (2012). Algoritmos: teoria e prática. Retrieved from \url{https://books.google.com.br/books?id=6iA4LgEACAAJ}.
        
        \bibitem[A04]{A04} Ascenio, Ana Fernanda Gomes. Estrutura de dados: Algoritmos, análise da complexidade e implementações em Java e C++. São Paulo: Pearson Prentice Hall, 2010.
        
        \bibitem[A05]{A05} Donald E. Knuth. Sorting and Searching, volume 3 of The Art ofComputer Programming. Addison-Wesley, 1973. Second edition, 1998
        
        \bibitem[A06]{A06} Szwarcfiter, Jayme Luiz. Estruturas de dados e seus algoritmos / Jayme Luiz Szwarcfiter, Lilian Markenzon. 3.ed. [Reimpr.]. - Rio de Janeiro : LTC, 2015.
        
\end{thebibliography}
\end{frame}
\section{}

%remover logo
\setbeamertemplate{logo}{}

\begin{frame}
\titlepage
\end{frame}

\end{document}
